Let $ L $ be a geometric lattice.
An element $ X $ is called \textit{modular} if
\begin{gather*}
r(X) + r(Y) = r(X \wedge Y) + r(X \vee Y) \quad \text{ for all } Y \in L,
\end{gather*}
where $ r $ denotes the rank function of $ L $.
A flat $ X $ of a simple matroid $ M $ is called \textit{modular} if $ X $ is modular in $ L(M) $, the lattice of flats of~$ M $.

\begin{Definition}
A simple matroid $ M $ on the ground set $ E $ is said to be a \textit{modular join} if there exist two proper modular flats $ E_{1} $ and $ E_{2} $ of $ M $ such that $ E = E_{1} \cup E_{2} $.
We also say that $ M $ is the \textit{modular join over $ X $}, denoted $ M = P_{X}(E_{1}, E_{2}) = P_{X}(M_{1}, M_{2}) $, where $ M_{i} \coloneqq M|E_{i} $ for $ i = 1,2 $ and $ X \coloneqq E_{1} \cap E_{2} $.
\end{Definition}

\begin{Remark}
Ziegler \cite{ziegler1991binary-potams} introduced a modular join, which is a special case of a \textit{generalized parallel connection} or a \textit{strong join} investigated by Brylawski~\cite{brylawski1975modular-totams} and Lindstr{\"o}m~\cite{lindstrom1978strong-joctsa}.
Our definition of a modular join is different from Ziegler's one.
However, they are equivalent (see \cite[Propositioin~3.3]{ziegler1991binary-potams} and \cite[Proposition~5.10]{brylawski1975modular-totams}).
In addition, note that~$ X $ is modular in~$ M_{1}$,~$M_{2} $, and~$ M $.
\end{Remark}

\begin{Definition}A matroid $ M $ is called \textit{round} (or \textit{nonsplit}) if the ground set is not the union of two proper flats. A subset $ S $ of the ground set of $ M $ is \textit{round} if the restriction $ M|S $ is round.
\end{Definition}

It is well known that the graphic matroid of a simple graph without isolated vertices is round if and only if the graph is complete (see \cite[Theorem~4.2]{borissova2016regular-etpad} for example) and that any induced subgraph isomorphic to a complete graph corresponds to a modular flat of a graphic matroid (see \cite[Proposition~6.9.11 and below]{oxley2011matroid} for example).
Our generalization of chordal graphs is defined as follows.
\begin{Definition}\label{modularly extended matroid}
Let $ \mathcal{ME} $ be the minimal class of simple matroids which satisfies the following conditions.
\begin{enumerate}[$(i)$]\itemsep=0pt
\item The empty matroid is a member of $ \mathcal{ME} $.
\item If a simple matroid $ M $ has a modular coatom $ X $ and $ M|X \in \mathcal{ME} $, then $ M \in \mathcal{ME} $.
\item Let $ M $ be a modular join of $ M_{1} $ and $ M_{2} $ over a round flat.
If $ M_{1}, M_{2} \in \mathcal{ME} $, then $ M \in \mathcal{ME} $.
\end{enumerate}
We say that a simple matroid in $ \mathcal{ME} $ is \textit{modularly extended}.
\end{Definition}

Comparing the conditions in Theorem \ref{Dirac} and Definition \ref{modularly extended matroid}, we can prove that a simple graphic matroid is modularly extended if and only if the associated graph is chordal.

Recent studies \cite{suyama2019signed-dm} and \cite{torielli2020freeness-a} treat a similar class for signed graphs and their associated arrangements.
Moreover, we will study modularly extended matroids associated with gain graphs in Section~\ref{sec:applications}.


The linear dependence of an arrangement $ \mathcal{A} $ determines a simple matroid $ M(\mathcal{A}) $ on itself.
Namely, a subset $ \{H_{1}, \dots, H_{n}\} \subseteq \mathcal{A} $ is defined to be independent if the codimension of the intersection $ H_{1} \cap \dots \cap H_{n} $ is equal to $ n $.
The main theorem of this paper is as follows (see Definitions~\ref{definition divisional flag} and~\ref{definition divisional freeness} for the definition of divisional flags and divisional freeness).
\begin{Theorem}\label{main divisionally free}
Every modularly extended matroid has a divisional flag. In particular, if the linear dependence matroid $ M(\mathcal{A}) $ on an arrangement $ \mathcal{A} $ is modularly extended, then $ \mathcal{A} $ is divisionally free.
\end{Theorem}
